\section{Background}

Particle collision detection is a fundamental component of many computational methods involving discrete entities, including granular dynamics, molecular simulation, rigid-body systems, multiphase flow, and particle-based continuum methods. In such systems, the identification of interacting neighbors often dominates total computational cost, particularly as particle count increases. Efficient collision-detection strategies are therefore essential for large-scale simulation.

The most direct approach is the all-pairs search, in which every particle is tested against every other particle. For a system of $N$ particles, this requires

\[
\mathcal{O}(N^2)
\]

pairwise distance evaluations and rapidly becomes impractical for large particle populations.

To reduce this cost, spatial acceleration structures are commonly employed. Uniform grids, linked-cell lists, spatial hashing, octrees, bounding-volume hierarchies, and sweep-and-prune methods partition space so that only nearby particles are considered as collision candidates. For sufficiently uniform particle distributions, these approaches may reduce expected cost toward approximately linear complexity.

With the emergence of graphics processing units (GPUs), collision detection has increasingly been implemented using massively parallel architectures. Most modern GPU methods perform broad-phase processing through compute shaders or general-purpose GPU kernels. Common strategies include parallel binning, radix sorting of cell keys, prefix-sum compaction, neighbor-list generation, and atomic insertion into shared cell structures. Although effective, these methods often incur synchronization overhead, scattered memory access, or preprocessing passes that compete with simulation workloads.

The graphical Particle Collision Detection (gPCD) method is motivated by the observation that particle occupancy generation can be expressed as a geometric placement problem rather than solely as a compute-sorting problem. By representing each particle with an axis-aligned bounding region (AARB) and assigning occupancy through corner placement during rasterization, broad-phase candidate generation can occur concurrently with normal graphics-pipeline execution.

This approach creates a hybrid graphics--compute workflow in which the rasterization stages construct the potentially colliding set, while compute shaders perform exact narrow-phase collision testing. Such a division seeks to improve hardware utilization, reduce preprocessing overhead, and provide scalable collision detection for large particle populations.

